% GENERATED FILE. Edit canonical lessons, then run npm run generate:books.
% canonical-source-hash: fnv1a64:a119ff12eac2377a
% canonical-lessons: GU-C177-varavar, GU-C177-avarnavar, GU-C177-vakhat, GU-C177-niyamit, GU-C177-rojindu

\chapter{Again and again}
\label{ch:gu-again-and-again}
\hlchaptermodality{177}

\begin{chapteropening}
\textbf{By the end of this chapter:} \emph{I can say how often I will do something.}

Complete the last lesson of chapter 177: I can say how often I will do something.
\end{chapteropening}

\section[\texorpdfstring{\gu{વારંવાર} (vārãvār)}{વારંવાર (vārãvār)}]{\gu{વારંવાર} (vārãvār) — again and again}

\label{lesson:GU-C177-varavar}



\begin{quote}
Before the new one: say the Gujarati for three quarters, then the Gujarati for almost.
\end{quote}

\subsection*{You'll want to know: \gu{વારંવાર}}
\textbf{\gu{વારંવાર}} — \emph{vārãvār} — \textquotedblleft{}again and again\textquotedblright{}.

\begin{quote}
\textbf{\gu{તે} \gu{વારંવાર} \gu{ફોન} \gu{કરશે}.} — \emph{te vārãvār phon karśe} — \textquotedblleft{}She will phone again and again.\textquotedblright{}
\end{quote}

\begin{cousinweb}
\textbf{\gu{વારંવાર}} (\emph{vārãvār}) is built by doubling \emph{vār}, \textquotedblleft{}a turn, a time\textquotedblright{}, which CDIAL traces to Sanskrit \emph{vāra}, \textquotedblleft{}appointed time, one's turn; day of the week\textquotedblright{} (CDIAL 11547).
\end{cousinweb}

\begin{grammarlens}[title={What you've built}]
One more word for everyday life at A2.
\end{grammarlens}

\subsection*{Your turn}
\begin{itemize}
\raggedright
  \item \emph{Say it:} \emph{vārãvār}
  \item \emph{Say it:} \emph{vārãvār}, once more
  \item \emph{Say it:} \emph{te vārãvār phon karśe} — She will phone again and again.
  \item \emph{Recall:} say the Gujarati for three quarters, then the Gujarati for almost, then say \emph{vārãvār} again
\end{itemize}

\begin{tcolorbox}[breakable,colback=teal!4,colframe=teal!35!black,title={Before you move on}]
What does \gu{વારંવાર} mean? (\textquotedblleft{}Again and again\textquotedblright{}.) Say it once more.
\end{tcolorbox}
\section[\texorpdfstring{\gu{અવારનવાર} (avārnavār)}{અવારનવાર (avārnavār)}]{\gu{અવારનવાર} (avārnavār) — now and then}

\label{lesson:GU-C177-avarnavar}



\begin{quote}
Before the new one: say the Gujarati for almost, then the Gujarati for again and again.
\end{quote}

\subsection*{You'll want to know: \gu{અવારનવાર}}
\textbf{\gu{અવારનવાર}} — \emph{avārnavār} — \textquotedblleft{}now and then\textquotedblright{}.

\begin{quote}
\textbf{\gu{અમે} \gu{અવારનવાર} \gu{મળીશું}.} — \emph{ame avārnavār maḷīśũ} — \textquotedblleft{}We will meet now and then.\textquotedblright{}
\end{quote}

\begin{cousinweb}
\textbf{\gu{અવારનવાર}} (\emph{avārnavār}) is a rhyming compound built on \emph{vār}, \textquotedblleft{}a turn, a time\textquotedblright{}, from Sanskrit \emph{vāra} (CDIAL 11547); how its first part was formed is not settled.
\end{cousinweb}

\begin{grammarlens}[title={What you've built}]
One more word for everyday life at A2.
\end{grammarlens}

\subsection*{Your turn}
\begin{itemize}
\raggedright
  \item \emph{Say it:} \emph{avārnavār}
  \item \emph{Say it:} \emph{avārnavār}, once more
  \item \emph{Say it:} \emph{ame avārnavār maḷīśũ} — We will meet now and then.
  \item \emph{Recall:} say the Gujarati for almost, then the Gujarati for again and again, then say \emph{avārnavār} again
\end{itemize}

\begin{tcolorbox}[breakable,colback=teal!4,colframe=teal!35!black,title={Before you move on}]
What does \gu{અવારનવાર} mean? (\textquotedblleft{}Now and then\textquotedblright{}.) Say it once more.
\end{tcolorbox}
\section[\texorpdfstring{\gu{વખત} (vakhat)}{વખત (vakhat)}]{\gu{વખત} (vakhat) — occasion, a time}

\label{lesson:GU-C177-vakhat}



\begin{quote}
Before the new one: say the Gujarati for again and again, then the Gujarati for now and then.
\end{quote}

\subsection*{You'll want to know: \gu{વખત}}
\textbf{\gu{વખત}} — \emph{vakhat} — \textquotedblleft{}occasion, a time\textquotedblright{}.

\begin{quote}
\textbf{\gu{હું} \gu{બે} \gu{વખત} \gu{આવીશ}.} — \emph{hũ be vakhat āvīś} — \textquotedblleft{}I will come twice.\textquotedblright{}
\end{quote}

\textbf{\gu{બે} \gu{વખત}}, twice.

\begin{cousinweb}
\textbf{\gu{વખત}} (\emph{vakhat}) is a loan from Arabic \emph{waqt}, \textquotedblleft{}a time, an appointed time\textquotedblright{}, which reached Gujarati through Persian; EtymDB lists the Persian word. Its Arabic root is \emph{w-q-t}, \textquotedblleft{}to fix a time\textquotedblright{} (Lane).
\end{cousinweb}

\begin{grammarlens}[title={What you've built}]
One more word for everyday life at A2.
\end{grammarlens}

\subsection*{Your turn}
\begin{itemize}
\raggedright
  \item \emph{Say it:} \emph{vakhat}
  \item \emph{Say it:} \emph{vakhat}, once more
  \item \emph{Say it:} \emph{hũ be vakhat āvīś} — I will come twice.
  \item \emph{Recall:} say the Gujarati for again and again, then the Gujarati for now and then, then say \emph{vakhat} again
\end{itemize}

\begin{tcolorbox}[breakable,colback=teal!4,colframe=teal!35!black,title={Before you move on}]
What does \gu{વખત} mean? (\textquotedblleft{}Occasion, a time\textquotedblright{}.) Say it once more.
\end{tcolorbox}
\section[\texorpdfstring{\gu{નિયમિત} (niyamit)}{નિયમિત (niyamit)}]{\gu{નિયમિત} (niyamit) — regular, regularly}

\label{lesson:GU-C177-niyamit}



\begin{quote}
Before the new one: say the Gujarati for now and then, then the Gujarati for occasion.
\end{quote}

\subsection*{You'll want to know: \gu{નિયમિત}}
\textbf{\gu{નિયમિત}} — \emph{niyamit} — \textquotedblleft{}regular, regularly\textquotedblright{}.

\begin{quote}
\textbf{\gu{હવે} \gu{હું} \gu{નિયમિત} \gu{ચાલીશ}.} — \emph{have hũ niyamit cālīś} — \textquotedblleft{}From now on I will walk regularly.\textquotedblright{}
\end{quote}

\begin{cousinweb}
\textbf{\gu{નિયમિત}} (\emph{niyamit}) is Sanskrit \emph{niyamita}, \textquotedblleft{}restrained, regulated\textquotedblright{}, taken over whole as a learned loan, a tatsama, built on \emph{niyama}, \textquotedblleft{}restraint, rule\textquotedblright{}.
\end{cousinweb}

\begin{grammarlens}[title={What you've built}]
One more word for everyday life at A2.
\end{grammarlens}

\subsection*{Your turn}
\begin{itemize}
\raggedright
  \item \emph{Say it:} \emph{niyamit}
  \item \emph{Say it:} \emph{niyamit}, once more
  \item \emph{Say it:} \emph{have hũ niyamit cālīś} — From now on I will walk regularly.
  \item \emph{Recall:} say the Gujarati for now and then, then the Gujarati for occasion, then say \emph{niyamit} again
\end{itemize}

\begin{tcolorbox}[breakable,colback=teal!4,colframe=teal!35!black,title={Before you move on}]
What does \gu{નિયમિત} mean? (\textquotedblleft{}Regular, regularly\textquotedblright{}.) Say it once more.
\end{tcolorbox}
\section[\texorpdfstring{\gu{રોજિંદું} (rojindũ)}{રોજિંદું (rojindũ)}]{\gu{રોજિંદું} (rojindũ) — everyday, daily}

\label{lesson:GU-C177-rojindu}



\begin{quote}
Before the new one: say the Gujarati for occasion, then the Gujarati for regular.
\end{quote}

\subsection*{You'll want to know: \gu{રોજિંદું}}
\textbf{\gu{રોજિંદું}} — \emph{rojindũ} — \textquotedblleft{}everyday, daily\textquotedblright{}.

\begin{quote}
\textbf{\gu{હું} \gu{રોજિંદું} \gu{કામ} \gu{સવારે} \gu{કરીશ}.} — \emph{hũ rojindũ kām savāre karīś} — \textquotedblleft{}I will do the everyday work in the morning.\textquotedblright{}
\end{quote}

Built on \textbf{\gu{રોજ}}, \emph{roj}, every day.

\begin{cousinweb}
\textbf{\gu{રોજિંદું}} (\emph{rojindũ}) is built on the taught \emph{roj}, \textquotedblleft{}every day\textquotedblright{}, which is usually traced to Persian \emph{rōz}, \textquotedblleft{}day\textquotedblright{}.
\end{cousinweb}

\begin{grammarlens}[title={What you've built}]
One more word for everyday life at A2.
\end{grammarlens}

\subsection*{Your turn}
\begin{itemize}
\raggedright
  \item \emph{Say it:} \emph{rojindũ}
  \item \emph{Say it:} \emph{rojindũ}, once more
  \item \emph{Say it:} \emph{hũ rojindũ kām savāre karīś} — I will do the everyday work in the morning.
  \item \emph{Recall:} say the Gujarati for occasion, then the Gujarati for regular, then say \emph{rojindũ} again
\end{itemize}

\begin{tcolorbox}[breakable,colback=teal!4,colframe=teal!35!black,title={Before you move on}]
What does \gu{રોજિંદું} mean? (\textquotedblleft{}Everyday, daily\textquotedblright{}.) Say it once more.
\end{tcolorbox}
